\documentclass[tikz,border=5pt]{standalone}
\usepackage[T1]{fontenc}
\usepackage{lmodern}
\usetikzlibrary{backgrounds}

\begin{document}
% Set this before including the data to show an intermediate greedy state.
% At step k, accepted merge edges 1,...,k are highlighted. The default is the
% final state. Keeping this figure-specific default in the body makes it
% available when standalone embeds the figure source into the main thesis.
\providecommand{\CreditGroupingStep}{7}

% Keep figure-specific colors in the document body so they are available both
% when this file is compiled independently and when standalone embeds its body
% into the main thesis after an external build is unavailable.
\definecolor{creditGroupOne}{HTML}{E15759}
\definecolor{creditGroupTwo}{HTML}{F28E2B}
\definecolor{creditGroupThree}{HTML}{76B7B2}
\definecolor{creditGroupFour}{HTML}{59A14F}
\definecolor{creditGroupFive}{HTML}{4E79A7}
\definecolor{creditAccepted}{HTML}{17823B}

\begin{tikzpicture}[
  x=1cm,
  y=1cm,
  credit node/.style={
    draw=black!65,
    rounded corners=2pt,
    line width=0.55pt,
    minimum height=8mm,
    text width=19mm,
    align=center,
    inner sep=2.5pt,
    font=\sffamily\scriptsize\bfseries,
    text=black
  },
  retained edge/.style={draw=black!30, line width=0.55pt},
  accepted edge/.style={draw=creditAccepted, line width=2.1pt},
  edge value/.style={
    midway,
    sloped,
    allow upside down,
    fill=white,
    fill opacity=0.88,
    text opacity=1,
    inner sep=0.7pt,
    font=\sffamily\tiny,
    text=black!75
  }
]

% The data records call these two renderer commands. Styling and graph data
% therefore remain independent and can be overridden in later figures.
\newcommand{\CreditNode}[5]{%
  \expandafter\gdef\csname CreditNodeLabel@#1\endcsname{#2}%
  \node[credit node, fill=creditGroup#5!72] (#1) at (#3,#4) {#2};%
}

\newcommand{\CreditEdge}[5]{%
  \ifnum#4=0
    \draw[retained edge] (#1) -- (#2)
      node[edge value, pos=#5] {#3};%
  \else
    \ifnum#4>\CreditGroupingStep
      \draw[retained edge] (#1) -- (#2)
        node[edge value, pos=#5] {#3};%
    \else
      \draw[accepted edge] (#1) -- (#2)
        node[edge value, pos=#5, text=creditAccepted!70!black,
             font=\sffamily\tiny\bfseries] {#3};%
    \fi
  \fi
}

% Create nodes first, then draw all edges on the background layer. Loading the
% records twice keeps a single source of truth while ensuring that crossing
% edges never obscure node labels.
\let\CreditNodeRenderer\CreditNode
\let\CreditEdgeRenderer\CreditEdge
\renewcommand{\CreditEdge}[5]{}%
% Data for the New German Credit association graph.
%
% Source: FELICS data/new_german_credit/data.csv, processed as in
% src/intervention_generation/tabular_intervention_generator.py with
% p < 0.01, Cramer's V >= 0.5, and n_max = 3. Values are rounded only for
% display. The final argument of \CreditEdge is the greedy merge step; zero
% marks a retained association that was not selected for the spanning forest.
%
% Keeping the records separate from the renderer allows the same graph to be
% reused with another \CreditGroupingStep or alternative edge styles.

% \CreditNode{id}{label}{x}{y}{final group}
\CreditNode{income}{Household Income}{-0.8}{4.0}{One}
\CreditNode{liquid}{Liquid Assets}{2.7}{3.4}{One}
\CreditNode{schufa}{SCHUFA Score (\%)}{1.1}{-2.4}{One}

\CreditNode{amount}{Credit Amount}{3.9}{0.9}{Two}
\CreditNode{purpose}{Credit Purpose}{1.5}{-0.2}{Two}
\CreditNode{term}{Credit Term (Months)}{3.8}{-3.8}{Two}

\CreditNode{realestate}{Real Estate Value}{-0.9}{0.8}{Three}
\CreditNode{ownhouse}{Lives in Own House}{3.0}{-1.4}{Three}
\CreditNode{age}{Applicant Age}{-2.1}{-3.8}{Three}

\CreditNode{job}{Job Category}{-4.0}{2.4}{Four}
\CreditNode{debt}{Monthly Debt Burden}{-5.0}{-1.3}{Four}

\CreditNode{householdsize}{Household Size}{6.2}{1.2}{Five}

% \CreditEdge{first node}{second node}{Cramer's V}{merge step}{label position}
\CreditEdge{realestate}{ownhouse}{0.90}{1}{0.55}
\CreditEdge{income}{liquid}{0.71}{2}{0.54}
\CreditEdge{amount}{term}{0.68}{3}{0.55}
\CreditEdge{amount}{purpose}{0.65}{4}{0.48}
\CreditEdge{debt}{job}{0.61}{5}{0.48}
\CreditEdge{schufa}{liquid}{0.60}{6}{0.48}
\CreditEdge{age}{ownhouse}{0.59}{7}{0.55}

\CreditEdge{term}{purpose}{0.64}{0}{0.46}
\CreditEdge{income}{amount}{0.61}{0}{0.55}
\CreditEdge{income}{job}{0.59}{0}{0.46}
\CreditEdge{purpose}{job}{0.58}{0}{0.52}
\CreditEdge{job}{ownhouse}{0.57}{0}{0.52}
\CreditEdge{purpose}{realestate}{0.57}{0}{0.47}
\CreditEdge{liquid}{realestate}{0.56}{0}{0.48}
\CreditEdge{income}{purpose}{0.56}{0}{0.47}
\CreditEdge{age}{purpose}{0.56}{0}{0.52}
\CreditEdge{amount}{liquid}{0.55}{0}{0.43}
\CreditEdge{age}{realestate}{0.55}{0}{0.52}
\CreditEdge{purpose}{liquid}{0.55}{0}{0.45}
\CreditEdge{job}{realestate}{0.54}{0}{0.54}
\CreditEdge{job}{liquid}{0.53}{0}{0.48}
\CreditEdge{purpose}{ownhouse}{0.52}{0}{0.53}
\CreditEdge{income}{realestate}{0.52}{0}{0.52}
\CreditEdge{income}{schufa}{0.50}{0}{0.58}

\begin{scope}[on background layer]
  \renewcommand{\CreditNode}[5]{}%
  \let\CreditEdge\CreditEdgeRenderer
  % Data for the New German Credit association graph.
%
% Source: FELICS data/new_german_credit/data.csv, processed as in
% src/intervention_generation/tabular_intervention_generator.py with
% p < 0.01, Cramer's V >= 0.5, and n_max = 3. Values are rounded only for
% display. The final argument of \CreditEdge is the greedy merge step; zero
% marks a retained association that was not selected for the spanning forest.
%
% Keeping the records separate from the renderer allows the same graph to be
% reused with another \CreditGroupingStep or alternative edge styles.

% \CreditNode{id}{label}{x}{y}{final group}
\CreditNode{income}{Household Income}{-0.8}{4.0}{One}
\CreditNode{liquid}{Liquid Assets}{2.7}{3.4}{One}
\CreditNode{schufa}{SCHUFA Score (\%)}{1.1}{-2.4}{One}

\CreditNode{amount}{Credit Amount}{3.9}{0.9}{Two}
\CreditNode{purpose}{Credit Purpose}{1.5}{-0.2}{Two}
\CreditNode{term}{Credit Term (Months)}{3.8}{-3.8}{Two}

\CreditNode{realestate}{Real Estate Value}{-0.9}{0.8}{Three}
\CreditNode{ownhouse}{Lives in Own House}{3.0}{-1.4}{Three}
\CreditNode{age}{Applicant Age}{-2.1}{-3.8}{Three}

\CreditNode{job}{Job Category}{-4.0}{2.4}{Four}
\CreditNode{debt}{Monthly Debt Burden}{-5.0}{-1.3}{Four}

\CreditNode{householdsize}{Household Size}{6.2}{1.2}{Five}

% \CreditEdge{first node}{second node}{Cramer's V}{merge step}{label position}
\CreditEdge{realestate}{ownhouse}{0.90}{1}{0.55}
\CreditEdge{income}{liquid}{0.71}{2}{0.54}
\CreditEdge{amount}{term}{0.68}{3}{0.55}
\CreditEdge{amount}{purpose}{0.65}{4}{0.48}
\CreditEdge{debt}{job}{0.61}{5}{0.48}
\CreditEdge{schufa}{liquid}{0.60}{6}{0.48}
\CreditEdge{age}{ownhouse}{0.59}{7}{0.55}

\CreditEdge{term}{purpose}{0.64}{0}{0.46}
\CreditEdge{income}{amount}{0.61}{0}{0.55}
\CreditEdge{income}{job}{0.59}{0}{0.46}
\CreditEdge{purpose}{job}{0.58}{0}{0.52}
\CreditEdge{job}{ownhouse}{0.57}{0}{0.52}
\CreditEdge{purpose}{realestate}{0.57}{0}{0.47}
\CreditEdge{liquid}{realestate}{0.56}{0}{0.48}
\CreditEdge{income}{purpose}{0.56}{0}{0.47}
\CreditEdge{age}{purpose}{0.56}{0}{0.52}
\CreditEdge{amount}{liquid}{0.55}{0}{0.43}
\CreditEdge{age}{realestate}{0.55}{0}{0.52}
\CreditEdge{purpose}{liquid}{0.55}{0}{0.45}
\CreditEdge{job}{realestate}{0.54}{0}{0.54}
\CreditEdge{job}{liquid}{0.53}{0}{0.48}
\CreditEdge{purpose}{ownhouse}{0.52}{0}{0.53}
\CreditEdge{income}{realestate}{0.52}{0}{0.52}
\CreditEdge{income}{schufa}{0.50}{0}{0.58}

\end{scope}

% Compact legend. The node swatches denote membership in the final partition,
% not an additional statistic.
\node[anchor=west, font=\sffamily\scriptsize] at (-6.0,-5.0) {Final groups:};
\foreach \x/\group/\label in {
  -4.05/One/$s_1$,
  -3.10/Two/$s_2$,
  -2.15/Three/$s_3$,
  -1.20/Four/$s_4$,
  -0.25/Five/$s_5$
}{
  \node[draw=black!55, fill=creditGroup\group!72, rounded corners=1pt,
        minimum width=5mm, minimum height=3mm, inner sep=0pt] at (\x,-5.0) {};
  \node[anchor=west, font=\sffamily\tiny] at (\x+0.32,-5.0) {\label};
}
\draw[accepted edge] (1.45,-5.0) -- (2.25,-5.0)
  node[anchor=west, font=\sffamily\scriptsize, text=black] at (2.42,-5.0)
  {accepted merge};
\draw[retained edge] (4.35,-5.0) -- (5.15,-5.0)
  node[anchor=west, font=\sffamily\scriptsize, text=black] at (5.32,-5.0)
  {other retained association};

\end{tikzpicture}
\end{document}
